\documentclass[11pt]{article}
\usepackage[margin=1in]{geometry}
\usepackage{amsmath,amssymb,amsthm,lmodern}
\usepackage[colorlinks=true,linkcolor=blue,urlcolor=blue]{hyperref}
\newtheorem{theorem}{Theorem}
\setlength{\emergencystretch}{2em}
\title{Inactive-coordinate rotations refute permutation rigidity}
\author{Myeongjin Kang}
\date{25 October 2026}
\begin{document}
\maketitle
\begin{abstract}
A rational orthogonal rotation in the kernel coordinates of a rectangular diagonal matrix leaves the source's stationary expression diagonal while failing to be a permutation. We verify every condition, including the actual negative fractional matrix power, in Lean. This is a smaller instance of the inactive-block mechanism already posted by Shivam Patel, not a new-discovery claim.
\end{abstract}
\section{Source assertion and explicit data}
Conjecture 3 in Appendix D of Centorrino, Bullo and Russo \cite{source}, also MathDB 368725, asserts that orthogonal $U$ is a permutation whenever
\[
K=U^T L^T(LUDU^TL^T)^{-1/3}LU
\]
is diagonal, where $L\in\mathbb R^{m\times n}$ is rectangular diagonal and $D\in\mathbb R^{n\times n}$ is positive definite diagonal with distinct diagonal entries. The displayed power is the usual spectral power of the positive definite active block.

Take $m=1$, $n=3$, and
\[
 L=\begin{pmatrix}1&0&0\end{pmatrix},\qquad
 D=\operatorname{diag}(1,1/2,1/3),\qquad
 U=\begin{pmatrix}1&0&0\\0&3/5&-4/5\\0&4/5&3/5\end{pmatrix}.
\]
These diagonal entries are strictly decreasing, so the data also satisfy that stronger ordering convention. The active entry of $L$ is positive and its row rank is one.
\begin{theorem}
These matrices satisfy every stated hypothesis, $K$ is diagonal, and $U$ is not a permutation matrix or a signed permutation matrix.
\end{theorem}
\begin{proof}
The last two columns of $U$ have squared norm $(3/5)^2+(4/5)^2=1$ and inner product zero; they are perpendicular to its first column. Therefore $U^TU=UU^T=I_3$. The positive distinct entries of $D$ make it positive definite.

Since $LU=L$, direct multiplication gives
\[
LUDU^TL^T=[1].
\]
This block is positive definite, so its negative one-third power is well defined and equals $[1]$. Consequently
\[
 K=U^TL^TLU=L^TL=\operatorname{diag}(1,0,0).
\]
The entry $U_{22}=3/5$ is neither zero nor one, so $U$ is not a permutation matrix. It is also not a signed permutation, whose entries would be zero, one, or minus one.
\end{proof}
\section{Mechanism, prior status, and formal coverage}
The kernel coordinates of $L$ are invisible in $LU$. Rotating these coordinates changes $U$ without changing the active block or $K$. Distinct eigenvalues in $D$ do not recover that discarded information.

Patel's existing response to MathDB 368725 \cite{patel} exhibits this same mechanism in a $2\times4$ example and a general inactive-block family. We use a smaller $1\times3$ instance for a direct formal certificate. There is no claim of discovery, of the entire corrected rigidity classification, or of the source's separate trace-maximization theorem.

The file \nolinkurl{formal/PermutationRigidity.lean} defines all three matrices and the exact expression \nolinkurl{stationary}. It uses \nolinkurl{CFC.rpow}, the continuous-functional-calculus real matrix power with the standard positive-semidefinite order. The active-block identity and \nolinkurl{CFC.one_rpow} evaluate this actual operation; no fractional-power axiom is introduced. The matrix products, positive definiteness, distinct diagonal entries, rectangular diagonal form, diagonality of $K$, and exclusion of every mathlib permutation matrix are proved. The final theorem \nolinkurl{counterexample} bundles these assertions.

The release includes the independent review, dependency pins, and verifier. Lean 4.34.1 checks all local declarations with warnings treated as errors. The only permitted axioms are propositional extensionality, classical choice, and quotient soundness. The result is an exact rational matrix counterexample over the reals, not a numerical experiment or a perturbative argument.
\begin{thebibliography}{9}
\bibitem{source} V. Centorrino, F. Bullo, and G. Russo. arXiv:2506.06134v1 (2025), Appendix D, Conjecture 3. \url{https://arxiv.org/html/2506.06134v1}.
\bibitem{patel} Shivam Patel. Counterexample and corrected rigidity analysis, MathDB 368725, solution 7b45cbbd-c0dc-469c-9c7c-73272a398aef. \url{https://mathdb.com/p/368725}.
\end{thebibliography}
\end{document}
